% =====================================================================
%  COPY THIS WHOLE FILE INTO OVERLEAF (select all -> paste into main.tex)
%  Self-contained: 2-column, references embedded (no .bib, no BibTeX run).
%  Compiler: pdfLaTeX (Overleaf default) works as-is.
%
%  Ton's sections — §1 Introduction, §2 Problem Statement, §3 Related Works
%  Group: Flower — Explainable Reviewed-Item Retrieval  (OPTION C)
%  Author: Chidsanuphong "Ton" Pengchai (st127004), DSAI, AIT
%
%  NOTE: the Thai example query is given as an English gloss so this compiles
%  under pdfLaTeX. To show Thai script, switch the Overleaf compiler to
%  XeLaTeX and add in the preamble:
%     \usepackage{fontspec}\newfontfamily\thaifont{Noto Sans Thai}
%  then use {\thaifont อยากกินอะไรเปรี้ยวๆ เผ็ดๆ ไม่แพง}.
%
%  Numbers in §1.1 (Time Out ranking, street-food revenue, tourist spending,
%  Wongnai user counts) still need proper source citations before submission.
% =====================================================================
\documentclass[10pt,twocolumn]{article}

\usepackage[utf8]{inputenc}
\usepackage[T1]{fontenc}
\usepackage{amsmath,amssymb}
\usepackage{cite}
\usepackage[hidelinks]{hyperref}
\usepackage[margin=1in]{geometry}

\title{Explainable Reviewed-Item Retrieval:\\Searching Restaurants by Taste Preference}
\author{Chidsanuphong ``Ton'' Pengchai (st127004)\\Data Science and AI, AIT --- Group Flower}
\date{}

\begin{document}
\maketitle

\section{Introduction}

\subsection{Background of the Work}
Food occupies a central place in Thai culture and in Thailand's national economy,
and Thai cuisine---street food in particular---is treated as an instrument of
national soft power that drives tourism revenue. The global standing of this food
scene is measurable: Bangkok was ranked the world's No.~2 food city (and No.~1 in
Asia) in \emph{Time Out}'s \emph{Best Cities for Food 2026}, behind only Lima, with
an 80\% expert-panel score and 81\% of surveyed locals rating the city highly for
food quality and diversity. The street-food segment is not a marginal niche but the
backbone of the sector: Thailand has roughly 103{,}000 street-food operators, about
69\% of all food businesses, generating approximately 271{,}355 million baht---about
32.4\% of food-business revenue---making street food the single largest foodservice
segment. Food is also where visitors spend: food and beverage accounted for 23\% of
foreign tourists' spending in 2023 according to the Bank of Thailand's Q2~2024
Monetary Policy Report (second only to accommodation at 36\%). What distinguishes
Thai food choice from many other cuisines is the centrality of \emph{flavour
balance}---the interplay of spicy, sour, sweet, salty, and rich (fatty/umami)
notes---so that diners frequently choose a meal by desired taste profile rather than
by dish name.

The digital layer over this landscape is dominated by two incumbents. Wongnai (now
part of the LINE MAN Wongnai group) is Thailand's leading restaurant-review platform;
its app draws on user-generated reviews and feedback from over 5.7~million users,
with more than 2.4~million reviews and 84~million photos across more than one million
locations nationwide. Google Maps provides the other major review and discovery
surface. Both, however, expose discovery primarily through keyword search, category
filters, star ratings, and geographic proximity---none of which lets a user search by
the \emph{sensory experience} they want. Academic NLP has matured enough to make such
a system feasible: transformer sentence encoders~\cite{devlin2019bert,reimers2019sbert}
and multilingual models, together with Thai-specific pretrained models such as
WangchanBERTa~\cite{lowphansirikul2021wangchanberta} and the publicly released Wongnai
review corpus, provide the ingredients for processing review text---including
low-resource Thai---at scale.

\subsection{Why Conduct This Project}
Three gaps motivate the work. First, \textbf{there is no mainstream way to search
restaurants by taste preference}: a user who wants ``something sour and spicy that is
not expensive'' must translate that intent into keywords, guess at dish names, and
manually cross-read reviews. Second, the recommendation \textbf{cold-start problem} is
acute here---a tourist or newcomer has no interaction history and no ratings, so
collaborative-filtering systems that depend on behavioural logs~\cite{koren2009mf}
cannot serve them, and the natural-language, content-based paradigm is the appropriate
alternative~\cite{gao2021crs,kostric2024usage}. Third, the evidence needed to answer
such a query \textbf{already exists latent in review text}---diners constantly describe
spiciness, sourness, value, and atmosphere in their own words---yet this signal remains
untapped as a first-class retrieval mechanism, and prior work leaves open \emph{how}
that sentence-level evidence should be combined into a restaurant-level
decision~\cite{rird2023,zhang2020explainable}.

\subsection{Business and Research Understanding of the Problem}
The problem is felt by several groups. \textbf{Tourists and newcomers} bear a search
cost: they cannot express taste intent as keywords and cannot efficiently read long
review threads, leading to poor or risky choices. \textbf{Locals} face choice overload
across a very large restaurant catalogue. \textbf{Vendors}---especially small and
street-food operators---are poorly served by rating-and-proximity ranking, because a
distinctive but niche flavour offering is invisible to a user who cannot search for it.
Crucially, star ratings measure \emph{overall satisfaction}, not spicy/sour/sweet/rich
intensity, so the incumbents' interaction model captures none of this. The research
question is therefore not ``which restaurant is best'' but a \textbf{retrieval
question}: given a free-text preference query, can review text be used to retrieve and
rank restaurants in a way that is both accurate and self-justifying, and does the way
sentence-level evidence is \emph{fused} into a restaurant-level score materially affect
that accuracy? This reframing turns a subjective ``flavour recommender'' into a
measurable information-retrieval problem with human relevance judgements~\cite{rird2023}.

\subsection{Potential Impact and Benefits}
For \textbf{tourists and newcomers}, the system removes the language and knowledge
barrier by accepting intent in natural language and by returning, alongside each
recommendation, the exact review sentences that justify it---an explainable output that
builds trust~\cite{zhang2020explainable}. For \textbf{locals} it offers taste-driven
discovery beyond rating and distance. For \textbf{vendors} it creates a fairer surface
where a distinctive flavour profile can be matched to demand rather than buried beneath
aggregate stars. Methodologically, and in line with the scope fixed in
Section~\ref{sec:scope}, the project contributes: (i)~a \emph{systematic comparison of
sentence-to-restaurant fusion functions} (max, mean-of-top-$k$, rating-weighted sum),
evaluated on the labelled RIRD benchmark with nDCG@10, MAP, and Recall@10---evidence
currently missing because the anchor prior work reports only a single aggregation
method (top-$K$ mean) and only two metrics (R-Precision, MAP)~\cite{rird2023};
(ii)~an explainable, content-based retrieval pipeline that attaches justifying review
sentences to every ranked restaurant~\cite{zhang2020explainable,kostric2024usage};
(iii)~a \emph{cross-lingual demonstration} of the same language-agnostic pipeline on
the Thai Wongnai corpus, showing the method transfers to a low-resource language at the
review/sentence level; and (iv)~an optional \emph{flavour-feature ablation} that tests
whether explicit sensory features add value over pure dense retrieval, judged against
independent human labels rather than against its own domain assumptions.

\section{Problem Statement}

\subsection{Problem Definition}
We address a natural-language, review-grounded restaurant retrieval problem. Let
$R = \{r_1, \dots, r_n\}$ be a corpus of restaurants, where each restaurant $r_i$ is
associated with a set of reviews $D_i = \{d_{i,1}, \dots, d_{i,m}\}$, and each review is
decomposed into sentences $S_i = \{s_{i,1}, s_{i,2}, \dots\}$. Given a free-text
preference query $q$, the task is to produce a \textbf{ranking of the restaurants in
$R$ by relevance to $q$}. The system first scores each sentence against the query with a
sentence encoder, $\mathrm{score}(q, s_{i,j})$; it then applies a fusion function $f$
that aggregates the sentence-level scores of restaurant $r_i$ into a single
restaurant-level score
\begin{equation}
  \mathrm{Score}(q, r_i) = f\big(\{\,\mathrm{score}(q, s_{i,j}) : s_{i,j} \in S_i\,\}\big),
\end{equation}
used for ranking; the top-ranked restaurants are returned together with their
highest-scoring justifying sentences. The central research contribution concerns the
choice of $f$---whether \emph{max}, \emph{mean-of-top-$k$}, or \emph{rating-weighted
sum} yields the best ranking. This is the \emph{reviewed-item retrieval} formulation, in
which the two-level restaurant/review structure requires an explicit aggregation
step~\cite{rird2023}.

The formulation is language-agnostic. It is \emph{developed and quantitatively
evaluated} on the English RIRD benchmark~\cite{rird2023}, which supplies the restaurant
grouping, 100 preference queries, and human relevance labels needed to measure ranking
quality. The same pipeline is then \emph{demonstrated} on the Thai Wongnai corpus---for
example, on the query ``I want something sour and spicy that is not expensive''---at the
review/sentence level, illustrating cross-lingual transfer to a low-resource setting
where relevance labels are not yet available.

\subsection{Objectives}
\begin{enumerate}
  \item \textbf{Implement the retrieval pipeline.} Build sentence-level dense retrieval
        using transformer sentence embeddings~\cite{reimers2019sbert,karpukhin2020dpr}:
        split reviews into sentences, encode query and sentences into a shared vector
        space, and retrieve the top-$k$ sentences per query.
  \item \textbf{Implement and systematically compare fusion.} Compare fusion strategies
        (max, mean-of-top-$k$, rating-weighted sum) for aggregating sentence scores into
        a restaurant score---the primary research contribution, filling the gap left by
        the anchor work, which tested only the top-$K$ mean~\cite{rird2023}.
  \item \textbf{Benchmark against baselines on RIRD.} Compare the pipeline against
        random ranking, rating-only sorting, BM25 lexical
        matching~\cite{robertson2009bm25}, and embedding retrieval with mean fusion,
        reporting \textbf{nDCG@10, MAP, and Recall@10} on the RIRD relevance labels.
  \item \textbf{Deliver explainable output.} For every recommended restaurant, return the
        specific review sentences that justify its position in the
        ranking~\cite{zhang2020explainable,kostric2024usage}.
  \item \textbf{Demonstrate cross-lingual application (Thai).} Port the same pipeline to
        the Thai Wongnai corpus and show qualitative review/sentence-level retrieval for
        Thai preference queries. This is a \emph{demonstration}, not a measured claim:
        Wongnai lacks restaurant grouping, queries, and relevance labels, so no
        quantitative result is asserted on it.
  \item \textbf{Test the flavour-feature layer (optional ablation).} Evaluate whether
        adding explicit sensory-intensity features (built from the Thailand Foods
        gazetteer and Thai aspect signals) improves ranking over pure dense retrieval,
        judging the layer against \textbf{independent human labels} rather than against
        its own domain assumptions~\cite{pontiki2014semeval}.
\end{enumerate}

\subsection{Scope and Boundaries}
\label{sec:scope}
The project is deliberately scoped to \textbf{content- and review-based retrieval}. All
\emph{quantitative} claims are made on the labelled RIRD benchmark; the Thai Wongnai
corpus is used only for a \emph{qualitative} cross-lingual demonstration, and full Thai
\emph{restaurant-level} recommendation---which would require Wongnai metadata with
restaurant identifiers or a purpose-built labelled Thai query set---is left as future
work. The project does \textbf{not} attempt collaborative filtering; it does \textbf{not}
use click logs or any behavioural interaction data; and it does \textbf{not} build
per-user profiles. It is \textbf{single-turn} (one query in, one ranking out) rather
than a multi-turn conversational dialogue. It does not perform menu-price extraction,
delivery logistics, or real-time availability checking, and it does not \emph{generate}
new natural-language explanations---only \emph{select} existing review sentences as
evidence. Geographic and temporal personalisation are out of scope. These boundaries
keep the study focused on the measurable retrieval-and-fusion research question defined
in Section~2.1.

\section{Related Works}

\subsection{Reviewed-Item Retrieval and the Fusion Gap}
The closest prior work and the anchor for this project is the Reviewed-Item Retrieval
(RIR) task of Abdollah Pour et al.~\cite{rird2023}, which formalises exactly our
setting: given a natural-language query and items each described by many reviews, a
system must aggregate query--review similarities into an item-level score and rank
items. The authors distinguish \emph{early fusion} (average an item's reviews into one
embedding, then score) from \emph{late fusion} (score each review, then aggregate the
top-$K$), and release the RIRD dataset---50 Toronto restaurants, roughly 29{,}000
reviews, and 100 preference queries in five categories (indirect, negation, general,
detailed, contradictory) with binary relevance from five annotators. Their best
late-fusion contrastive model reaches MAP~0.626, well above BM25 and off-the-shelf
BERT. Crucially for us, the paper evaluates with \emph{only} R-Precision and MAP and
aggregates with \emph{only} a top-$K$ mean; it does not test max or rating-weighted
fusion, nor report nDCG or Recall. This is the precise, non-manufactured gap our fusion
comparison fills, on the same dataset and query protocol.

\subsection{From Collaborative Filtering to Cold-Start, Content-Based Recommendation}
The dominant classical recommendation paradigm is collaborative filtering via matrix
factorization~\cite{koren2009mf}, which learns latent user and item factors from a
history of ratings or interactions. This paradigm is powerful when behavioural logs are
abundant, but it is structurally unable to serve the cold-start user---a tourist or
newcomer with no history---which is precisely our target scenario. Surveys of
conversational and content-based recommendation~\cite{gao2021crs} argue that
natural-language interfaces let users express current, context-specific intent directly
and are especially valuable when interaction data is unavailable. Our system is a
single-turn instance of this natural-language paradigm: it deliberately avoids
collaborative filtering and instead retrieves over review content, which is exactly what
makes it applicable to users the classical paradigm cannot reach.

\subsection{Natural-Language Preference Elicitation at the Review-Sentence Level}
Kostric et al.~\cite{kostric2024usage} observe that ordinary users often cannot answer
questions about item \emph{attributes}---they know how they intend to use an item, not
its technical specifications---and show that review sentences are a mineable source of
user-facing, preference-relevant language. This supports two of our design decisions.
First, it validates operating at the \emph{review-sentence} granularity, the same level
at which we retrieve and return justifications. Second, its central premise matches our
tourist/newcomer use case exactly: a visitor cannot name a Thai dish but can type
``something sour and spicy, not expensive,'' which is why free-text taste queries are
preferable to fixed attribute filters.

\subsection{Lexical and Dense Text Retrieval}
Our retrieval stack builds on two families of methods. The lexical family is anchored by
BM25~\cite{robertson2009bm25}, the probabilistic relevance model that remains a strong
and standard baseline; we use it as the lexical baseline our embedding approach must
beat. The dense family begins with BERT~\cite{devlin2019bert} and its sentence-level
adaptation, Sentence-BERT~\cite{reimers2019sbert}, which produces semantically
meaningful sentence embeddings comparable by cosine similarity---the technique we use to
encode both queries and review sentences into a shared space. Dense Passage
Retrieval~\cite{karpukhin2020dpr} then established that dense embedding retrieval can
substantially outperform lexical matching for natural-language queries; our BM25-vs-dense
comparison on restaurant reviews is a structurally identical experiment applied to a new
domain.

\subsection{Explainable Recommendation}
The explainable component of our system---returning the specific review sentences that
justify each restaurant's ranking---sits squarely within the ``review-sentence-as-
explanation'' paradigm catalogued by Zhang and Chen's survey of explainable
recommendation~\cite{zhang2020explainable}. The survey organises the field by
information source and model type, and argues that user reviews are a rich explanation
source because individual sentences can both justify a recommendation and surface the
aspects a user cares about. It also stresses evaluating explanations against independent
signals rather than self-referentially---guidance we adopt for the optional
flavour-feature layer.

\subsection{Aspect-Based Sentiment and the Flavour-Feature Layer}
The optional flavour layer draws on aspect-based sentiment analysis (ABSA), defined as a
mainstream task by SemEval-2014~\cite{pontiki2014semeval}, whose restaurant-domain
taxonomy---\emph{food, service, price, ambience, anecdotes/miscellaneous}---the
multilingual M-ABSA resource (used in our project) inherits. This citation lets us state
a precise design boundary: the standard ABSA taxonomy contains \emph{no}
flavour-intensity axes such as spicy, sour, or sweet. Consequently, off-the-shelf ABSA
aspects cannot serve as ground truth for a flavour-intensity layer; the layer must
instead be validated against independent human labels. This both motivates the layer's
place as an optional ablation and directly answers the concern that ingredient- and
lexicon-derived flavour features might otherwise validate each other circularly.

\subsection{Thai Language Processing}
Because the cross-lingual demonstration operates on Thai text, which is written without
inter-word spaces and is comparatively low-resource, it relies on Thai-specific
language technology. WangchanBERTa~\cite{lowphansirikul2021wangchanberta} is a
transformer model pretrained on large Thai corpora and, together with multilingual
sentence encoders and the PyThaiNLP toolkit for word segmentation, makes review-level
retrieval over the Wongnai corpus feasible without English-centric assumptions.

% =====================================================================
%  REFERENCES (embedded — no .bib / BibTeX needed). IEEE-style numbering.
% =====================================================================
\begin{thebibliography}{99}

\bibitem{rird2023}
M.~M. Abdollah Pour, P.~Farinneya, A.~Toroghi, A.~Korikov, A.~Pesaranghader,
T.~Sajed, M.~Bharadwaj, B.~Mavrin, and S.~Sanner,
``Self-supervised contrastive BERT fine-tuning for fusion-based reviewed-item
retrieval,'' in \emph{Advances in Information Retrieval (ECIR 2023)}, LNCS vol.~13980.
Cham, Switzerland: Springer, 2023, pp.~3--17, doi: 10.1007/978-3-031-28244-7\_1.

\bibitem{koren2009mf}
Y.~Koren, R.~Bell, and C.~Volinsky, ``Matrix factorization techniques for recommender
systems,'' \emph{Computer}, vol.~42, no.~8, pp.~30--37, 2009,
doi: 10.1109/MC.2009.263.

\bibitem{devlin2019bert}
J.~Devlin, M.-W. Chang, K.~Lee, and K.~Toutanova, ``BERT: Pre-training of deep
bidirectional transformers for language understanding,'' in \emph{Proc. NAACL-HLT},
2019, doi: 10.18653/v1/N19-1423.

\bibitem{reimers2019sbert}
N.~Reimers and I.~Gurevych, ``Sentence-BERT: Sentence embeddings using Siamese
BERT-networks,'' in \emph{Proc. EMNLP-IJCNLP}, 2019, doi: 10.18653/v1/D19-1410.

\bibitem{karpukhin2020dpr}
V.~Karpukhin, B.~O\u{g}uz, S.~Min, P.~Lewis, L.~Wu, S.~Edunov, D.~Chen, and W.-t. Yih,
``Dense passage retrieval for open-domain question answering,'' in \emph{Proc. EMNLP},
2020, doi: 10.18653/v1/2020.emnlp-main.550.

\bibitem{robertson2009bm25}
S.~Robertson and H.~Zaragoza, ``The probabilistic relevance framework: BM25 and
beyond,'' \emph{Foundations and Trends in Information Retrieval}, vol.~3, no.~4,
pp.~333--389, 2009, doi: 10.1561/1500000019.

\bibitem{zhang2020explainable}
Y.~Zhang and X.~Chen, ``Explainable recommendation: A survey and new perspectives,''
\emph{Foundations and Trends in Information Retrieval}, vol.~14, no.~1, pp.~1--101,
2020, doi: 10.1561/1500000066.

\bibitem{gao2021crs}
C.~Gao, W.~Lei, X.~He, M.~de Rijke, and T.-S. Chua, ``Advances and challenges in
conversational recommender systems: A survey,'' \emph{AI Open}, vol.~2, pp.~100--126,
2021, doi: 10.1016/j.aiopen.2021.06.002.

\bibitem{kostric2024usage}
I.~Kostric, K.~Balog, and F.~Radlinski, ``Generating usage-related questions for
preference elicitation in conversational recommender systems,'' \emph{ACM Trans.
Recommender Systems}, vol.~2, no.~2, art.~12, pp.~1--24, 2024, doi: 10.1145/3629981.

\bibitem{pontiki2014semeval}
M.~Pontiki, D.~Galanis, J.~Pavlopoulos, H.~Papageorgiou, I.~Androutsopoulos, and
S.~Manandhar, ``SemEval-2014 Task~4: Aspect based sentiment analysis,'' in \emph{Proc.
8th Int. Workshop on Semantic Evaluation (SemEval 2014)}, Dublin, Ireland, 2014,
pp.~27--35, doi: 10.3115/v1/S14-2004.

\bibitem{lowphansirikul2021wangchanberta}
L.~Lowphansirikul, C.~Polpanumas, N.~Jantrakulchai, and S.~Nutanong,
``WangchanBERTa: Pretraining transformer-based Thai language models,''
arXiv:2101.09635, 2021.

\end{thebibliography}

\end{document}
